\subsection{The Comparison Under a Fixed Codec}
\label{sec:same_codec}

Tables~\ref{tab:storage_at_rmse}--\ref{tab:error_at_size} let the codec be
chosen independently for each region, each grid and each target. That
protocol answers a definite question, namely what the best of the tested
codec family can do on each representation, and the answer is an envelope,
since no archive selects its codec region by region after the fact. The
same question with the codec held fixed on both grids is the controlled,
operational form of the comparison, and it separates two things that the
envelope combines: whether a grid compresses better, and whether the codecs
that suit it best happen to differ from those that suit the other.
Table~\ref{tab:fixed_codec} reports it, under the same bracketing rule as
the oracle tables and with the number of qualifying regions in every cell.

One reporting rule differs here, deliberately. Elsewhere a target is reported
only when at least three quarters of the regions bracket it, because those
tables carry the paper's paired grid-level inference and a target established
on a selected minority would not support it. Table~\ref{tab:fixed_codec} is a
codec-by-codec sensitivity table rather than an inference, so every cell is
shown with its own $n$ and none is withheld: at the loosest target JPEG~XL
qualifies on 26 regions and JPEG~2000 on 30, below the three-quarter mark, and
those cells are descriptive. The verdicts of
Sections~\ref{sec:rd_main}--\ref{sec:rd_illustrative} rest on the withheld
rule; this table is read beside them, not in place of them.


Each cell is the median per-region HPX/UTM compressed-size ratio at that
target when the \emph{same} codec is used on both grids, the UTM side
carrying $\gamma_{\mathrm{tile}}$, with the number of qualifying regions in
parentheses. A region enters a cell only if that codec brackets the target
with two measured points on both grids.

\begin{table}[t]
\centering
\caption{Same-codec control}
\label{tab:fixed_codec}
\begin{tabular}{lrrrrr}
\toprule
Codec & $10^{-4}$ & $3\times10^{-4}$ & $10^{-3}$ & $3\times10^{-3}$ & $10^{-2}$ \\
\midrule
\texttt{zfp} & 1.92 & 1.66 & 1.61 & 1.64 & 1.75 \\
\texttt{sz3} & 1.85 & 1.97 & 1.43 & 1.33 & 1.32 \\
\texttt{jpeg2000\_lossy} & 1.86 & 1.72 & 1.73 & 1.50 & 1.69 \\
\texttt{jpegxl\_lossy} & 1.76 & 1.69 & 1.32 & 1.37 & 1.28 \\
\texttt{dwt\_zstd} & 1.89 & 1.63 & 1.57 & 1.51 & 1.16 \\
\bottomrule
\end{tabular}
\vspace{3pt}
\begin{minipage}{\columnwidth}
\footnotesize
Median per-region HPX/UTM compressed-size ratio at each target RMSE when the \emph{same} codec is used on both grids (the UTM side includes $\gamma_{\mathrm{tile}}$). It complements the oracle-best comparison of Table~\ref{tab:storage_at_rmse}, which lets codec selection differ between grids.
\end{minipage}
\end{table}


Two observations follow. The codec ranking and the grid comparison are
different results. JPEG~2000 and JPEG~XL lead on both grids alike
(Section~\ref{sec:codec_choice}), and at $10^{-2}$ the oracle selects one of
the two on 30 of the 32 qualifying regions on the HEALPix side and on 31 of
them on the native side, so the oracle ratios of
Table~\ref{tab:storage_at_rmse} are close to what either codec gives on its
own: $0.77$ and $0.71$ against $0.72$ for the envelope. The crossover is therefore not an artefact of codec selection.
It is, however, codec-dependent in a way that the envelope hides. Four of
the five codecs reproduce it, JPEG~XL and JPEG~2000 most strongly
($0.71$ and $0.77$ at $10^{-2}$, $p<10^{-7}$), then DWT+Zstandard
($0.81$) and SZ3 ($0.93$, $p=0.02$); ZFP does not reproduce it at all.
Under ZFP the HEALPix tile costs $1.54$--$1.72$ times the native product
at every target, with no decline as the target loosens. ZFP is a
fixed-accuracy coder whose block transform spends bits according to the
tolerance it is given rather than to the content it finds, and the
redundancy of an oversampled tile is precisely content that a
rate-adaptive coder can decline to spend bits on and a fixed-accuracy coder
cannot. That is an interpretation of a measured result, not a mechanism this
experiment establishes; SZ3, which also reproduces the crossover, is
prediction-based rather than a transform coder, so the successful family
cannot be defined by the transform alone. What is measured is that the
crossover is codec-dependent, reproduced by four of the five codecs tested
and not by ZFP. An archive that standardised on ZFP would see no regime in
which the carrier is the smaller representation; one that standardised on any
of the other four would see the regime of
Table~\ref{tab:storage_at_rmse}. The scope of the crossover is thus a carrier
populated by a transfer in the smoothness class of Clough--Tocher
(Section~\ref{sec:transfer_cost}) \emph{and} coded by one of the four, and the
second condition is as binding as the first.
